\textit{What the results mean.} A 30\,m elevation model keeps enough of the dune shape to rank face cells better than chance, but not enough to plan around faces: the best 30\,m layers add a few points of mission success, while the same layers on a 10\,m map add up to {{VAL:ho_post_E1_I10_BICUBICcal_d}} points. Fig.~\ref{fig:resolution}b suggests why: faces several grid cells long are found at every spacing, and most faces in these fields are shorter than one 30\,m cell. Better learning helps detection a little (super resolution {{VAL:auc_I30_SR}} against {{VAL:auc_I30_BICUBIC}} for interpolation) and planning about as much as interpolation (Section~\ref{sec:posthoc}). Imagery gave the largest gain in 30\,m detection, as in dune mapping~\cite{hugenholtz2012remote}, but not in planning. With a sensor that sees every face nearby the prior hardly matters, but faces hide behind crests, and with line of sight a good face map more than doubled success. The face map results put the remaining gap on information: face information from fine data, placed to within a few metres, would make most 30\,m routes safe, so a drone survey must be georeferenced accordingly and its face threshold calibrated for its own resolution.

\textit{Arabian dunes.} We did not test them. A face 100\,m high at 32\,degrees spans about 160\,m downslope, five 30\,m cells, so large megadune faces should be visible at 30\,m and smaller flank faces not; but GLO-30 is weaker over sand~\cite{martone2016enhancing} and dunes migrate.

\textit{Limitations.} The truth is steep ground: rock, vehicle ruts and vegetated mounds also exceed 28\,degrees, Dumont, Algodones and Little Sahara are off road vehicle areas, and the manual audit of a stratified sample of 600 patches is not finished, so label precision is unknown. We chose the windows, used them in earlier tests of the detector, and keep the prediction log in our own repository. The vehicle model is a posterior fitted to published slip curves of other vehicles; in our own leave one study out validation no slip model family transferred between published studies (median held out slip error 0.20 to 0.41 depending on the family), so absolute route outcomes are uncertain. No rover was driven, no physics was simulated, and even the line of sight check assumes that every visible face is detected. Six sites limit the training data of the learned layers and the power of site level tests.
